\chapter{Resolution of identity (RI/DF)}
\label{ch:ri-df}

The four-center two-electron integrals of the multireference chain are
the memory bottleneck of its nuclear-derivative path.  Under
\texttt{pyscfad}, the derivative of \texttt{int2e} allocates a dense
\texttt{nctr}~$\times$~\texttt{nao}$^4$ array before any chemistry
happens: 55.0~GB at \texttt{nao}~$=96$ and 2033~GB at \texttt{nao}~$=175$
for the systems measured here.  The module \texttt{diffmrc.core.ri}
replaces the four-center tensor in the correlation stage by two
three-center factors and a metric, and supplies a nuclear-derivative
path whose working set is \texttt{nao}$^2$~$\times$~\texttt{naux} per
atom.  This chapter states the approximation, the public interface, the
axis and sign conventions, the measured error budget, the memory
figures, and the boundaries of the implementation.

\section{The approximation}

The RI (density-fitting) approximation writes

\begin{equation}
(pq|rs) \;\approx\; \sum_{P Q} (pq|P)\, S^{-1}_{PQ}\, (Q|rs),
\qquad S_{PQ} = (P|Q),
\label{eq:ri}
\end{equation}

where capital indices run over an auxiliary basis and $S$ is the
auxiliary metric.  With the three-center tensor
$B_{\mu\nu P} = (\mu\nu|P)$ and the symmetric inverse square root
$J = S^{-1/2}$, the four-center tensor is assembled as

\begin{equation}
W = B J, \qquad
(\mu\nu|\lambda\sigma) = \sum_{P} W_{\mu\nu P} W_{\lambda\sigma P},
\label{eq:ri-assembly}
\end{equation}

that is, one matrix product followed by one GEMM over the pair index.
The two-step form is the implementation, not an optimization: an
algebraically equivalent single four-operand einsum showed absolute
errors of order $10^{-6}$ on a $96^4$ array on this stack, and two such
chains differed by eight percent in the rotation gradient.

The inverse square root is computed from the symmetric eigendecomposition
of $S$.  Directions with eigenvalue below \texttt{rcond}~$\times$~the
largest eigenvalue are truncated to zero in $J$; the default
\texttt{rcond} is $10^{-12}$, and the number of dropped directions is
reported by the factor container.

The derivative path uses the factor-level form of
Eq.~\eqref{eq:ri-assembly}: the host code supplies $B$, $J$ and their
nuclear derivatives, and a pure-JAX function with an explicit
\texttt{custom\_jvp} maps $(B, J)$ to the four-center tensor.  No
four-center derivative integral is ever formed, and no
eigendecomposition is differentiated inside JAX; the derivative of $J$
is the analytic Daleckii--Krein divided difference of $x^{-1/2}$, which
is finite at degenerate metric eigenvalues.

\section{API reference}

The functions below are exported by \texttt{core.ri}.  Shapes follow the
conventions of Section~\ref{sec:ri-conventions}.

\begin{center}
\begin{tabularx}{\textwidth}{>{\raggedright\arraybackslash}p{0.40\textwidth}>{\raggedright\arraybackslash}X}
\toprule
Entry point & Returns and purpose \\
\midrule
\texttt{make\_auxmol(mol, auxbasis=None)} &
  \texttt{(auxmol, label)}.  Tries the names in
  \texttt{AUX\_BASIS\_PRIORITY} and falls back to the orbital basis;
  a basis that produces zero functions counts as a failure. \\
\texttt{ri\_three\_center(mol, auxmol)} &
  $B$, shape \texttt{(nao, nao, naux)}; uses
  \texttt{df.incore.aux\_e2}, not \texttt{mol.intor('int3c2e\_sph')}. \\
\texttt{ri\_metric(auxmol)} &
  $S$, shape \texttt{(naux, naux)}, symmetrized. \\
\texttt{inv\_sqrt\_metric(S, rcond, eig=None)} &
  Truncated symmetric $S^{-1/2}$, shape \texttt{(naux, naux)}. \\
\texttt{ri\_factors(mol, auxmol, auxbasis=None, rcond)} &
  \texttt{RIFactors} carrying $B$, $S$, $J$, the metric
  eigendecomposition, the basis label and \texttt{rcond}. \\
\texttt{ri\_eri\_ao(mol, auxmol, factors=None, rcond)} &
  RI four-center tensor, shape \texttt{(nao, nao, nao, nao)}, chemist
  notation. \\
\texttt{ri\_eri\_mo(C, b, jhalf=None, full=True)} &
  \texttt{(b\_mo, eri\_mo)} with shapes
  \texttt{(norb, norb, naux)} and \texttt{(norb, norb, norb, norb)};
  \texttt{eri\_mo} is \texttt{None} when \texttt{full=False}. \\
\texttt{ri\_three\_center\_deriv(mol, auxmol, atom)} &
  $dB/dR_{\mathrm{atom}}$, shape \texttt{(3, nao, nao, naux)}. \\
\texttt{ri\_metric\_deriv(auxmol, atom)} &
  $dS/dR_{\mathrm{atom}}$, shape \texttt{(3, naux, naux)}. \\
\texttt{ri\_factors\_deriv(mol, auxmol, atom, factors)} &
  \texttt{RIFactorDerivs} with \texttt{db}, \texttt{ds} and
  \texttt{djhalf}. \\
\texttt{ri\_eri\_deriv\_contract(gamma, factors, derivs, per\_axis=False)} &
  Contracts $\gamma = \partial E/\partial(\text{eri\_ao})$ with one
  atom's derivative; scalar, or a length-3 vector when
  \texttt{per\_axis=True}. \\
\texttt{ri\_eri\_from\_factors(b, jhalf)} &
  Pure-JAX RI four-center tensor with the analytic \texttt{custom\_jvp}. \\
\texttt{ri\_eri\_jvp(b, jhalf, db, djhalf)} &
  Convenience wrapper returning the value and the directional
  derivative. \\
\bottomrule
\end{tabularx}
\end{center}

Two further entries exist for validation only.  \texttt{ri\_eri\_deriv}
materializes the full \texttt{(3, nao, nao, nao, nao)} derivative and
refuses to allocate above a caller-supplied byte limit; production code
uses the factor derivatives with \texttt{ri\_eri\_deriv\_contract} or
the JAX rule instead.  \texttt{inv\_sqrt\_metric\_jvp} exposes the
metric derivative on its own.

The intended nuclear-gradient recipe is:

\begin{lstlisting}[language=Python]
from diffmrc.core import ri

factors = ri.ri_factors(mol, auxmol)              # host, numpy
E = energy_from_factors(factors.b, factors.jhalf) # pure JAX chain
dE_db, dE_djhalf = jax.grad(E, argnums=(0, 1))(
    factors.b, factors.jhalf)
for atom in range(mol.natm):
    d = ri.ri_factors_deriv(mol, auxmol, atom, factors)
    dE_dR[atom] = (np.sum(dE_db * d.db)
                   + np.sum(dE_djhalf * d.djhalf))
\end{lstlisting}

The only derivative integrals touched are the three-center
\texttt{ip1}/\texttt{ip2} arrays and the metric \texttt{ip1}/\texttt{ip2}
arrays.

\section{Axis and sign conventions}
\label{sec:ri-conventions}

The conventions are fixed in the module and verified by isolated finite
differences; they are repeated here because the \texttt{libcint} names
are misleading.

\begin{itemize}
\item $B[\mu,\nu,P] = (\mu\nu|P)$: both orbital indices first, the
  auxiliary index last.  $S[P,Q] = (P|Q)$.
\item $\texttt{eri\_ao}[\mu,\nu,\lambda,\sigma] =
  (\mu\nu|\lambda\sigma)$ in chemist notation.
\item \texttt{int3c2e\_ip1} differentiates the \emph{first} orbital
  center and \texttt{int3c2e\_ip2} differentiates the \emph{auxiliary}
  center, not the second orbital center.  \texttt{int3c2e\_ip3} is not
  exported by \texttt{libcint} on this stack.
\item The \texttt{ip} series carries the electronic-coordinate sign:
  $\partial(\mu\nu|P)/\partial R_{c(\mu)} = -\texttt{ip1}$.  The second
  orbital center is $-\texttt{ip1}$ with the orbital-pair axes
  exchanged, and the auxiliary center is $-\texttt{ip2}$.
\item For the metric,
  $dS/dR_A = -(\texttt{ip1}[P\in A] + \texttt{ip2}[Q\in A])$.
\end{itemize}

The isolated checks displace the molecule with the auxiliary basis fixed
and the auxiliary basis with the molecule fixed: the two groups of terms
reproduce their finite differences to $7.9\times10^{-11}$, and the sum
reproduces the full displacement to $4.8\times10^{-11}$.

\section{Error budget}

Three error classes are reported separately, because they have
different origins and different remedies.

\begin{center}
\begin{tabularx}{\textwidth}{>{\raggedright\arraybackslash}p{0.18\textwidth}X}
\toprule
Class & Measured value and meaning \\
\midrule
Implementation &
  RI assembly against the pyscf DF reconstruction:
  $5.8\times10^{-13}$ (H$_2$O/6-31G) and $1.3\times10^{-12}$
  (N$_2$/cc-pVDZ).  This is the numerical floor of the contraction,
  not a physical error. \\
RI physical &
  The approximation of Eq.~\eqref{eq:ri} at fixed orbitals and CI
  vector: NEVPT2 $\Delta e_2$ of $+1.1\times10^{-5}$ (N$_2$/cc-pVDZ)
  and $-7.2\times10^{-6}$~Ha (CO/cc-pVDZ); $-1.8\times10^{-6}$~Ha on
  ThO/6-31G with the \texttt{autoaux} fitting set. \\
Chain derivative &
  RI-CASCI gradient against a finite difference of the same RI energy:
  $9.6\times10^{-8}$ (H$_2$O/6-31G) and $3.1\times10^{-7}$
  (N$_2$/cc-pVDZ), absolute, in Eh/bohr. \\
\bottomrule
\end{tabularx}
\end{center}

The derivative integrals themselves agree with central differences at a
$10^{-4}$~bohr step to $3.6\times10^{-9}$ (three-center tensor) and
$1.2\times10^{-9}$ (metric) on H$_2$O/6-31G.  The materialized ERI
derivative shows $2.5\times10^{-8}$ on the same system with the default
fitting basis, whose metric condition number is $2.5\times10^{7}$;
switching to a fitting basis with condition number $2.5\times10^{5}$
brings every component below $6.7\times10^{-9}$.  The residual on the
default basis is the finite-difference floor amplified by the metric
condition number, not an implementation defect.

The physical error is controlled by the fitting basis, not by the code.
The values above sit at the lower edge of the usual RI band
($10^{-5}$--$10^{-3}$~Ha for correlation energies) and must be quoted
with the fitting basis and the system attached.

\section{Memory}

The value path still materializes the dense \texttt{nao}$^4$ tensor,
because that is the object the CI layer consumes.  The derivative path
does not: its per-atom working set is three three-index arrays of size
\texttt{nao}$^2$~$\times$~\texttt{naux}, plus the metric derivative.

\begin{center}
\begin{tabular}{lrrrr}
\toprule
\texttt{nao} & \texttt{naux} & RI chain peak RSS & dense derivative array
  & array ratio \\
\midrule
96  & 646 & 2.11~GB  & 55.0~GB  & $26\times$ \\
175 & 746 & 16.05~GB & 2033~GB  & $127\times$ \\
\bottomrule
\end{tabular}
\end{center}

The peaks cover the integral-and-derivative chain only, without SCF or
CI.  The \texttt{nao}~$=96$ row uses the orbital-basis fallback and the
\texttt{nao}~$=175$ row the \texttt{autoaux} route; the memory scaling
depends on \texttt{naux}, not on which fitting route produced it.  The
derivative-array ratio in the table is $26\times$ at
\texttt{nao}~$=96$ and $127\times$ at \texttt{nao}~$=175$.  Against the
measured dense whole-chain peak of 210--213~GB at \texttt{nao}~$=96$,
the end-to-end reduction is about $100\times$.  The per-atom derivative
factors occupy 152.9~MB at \texttt{nao}~$=96$.

\section{Boundaries}

\begin{itemize}
\item \emph{The SCF layer remains dense.}  RI removes the memory wall
  of the correlation stage, not of the reference determinant.  A
  full-chain nuclear gradient for an actinide still requires an RI-JK
  SCF or an explicit frozen-reference partial-derivative semantics.
\item \emph{Second derivatives are not implemented.}  The factor-level
  \texttt{custom\_jvp} is first order.  Differentiating the reverse pass
  a second time is possible in principle but is outside the validated
  surface; the analytic $S^{-1/2}$ rule is available first-order only.
\item \emph{Actinide fitting sets are missing.}  The def2 and
  cc-pVDZ-RI fitting families stop at Rn.  For Th the module falls back
  to the orbital basis (\texttt{naux}~$=646$, metric condition number
  $1.5\times10^{9}$), which shifts a CASCI energy by $-0.594$~mHa.  The
  even-tempered construction \texttt{aug\_etb\_for\_dfbasis} is the best
  available substitute ($+0.021$~mHa, \texttt{naux}~$=625$); a dedicated
  actinide JKfit basis remains an open item.
\item \emph{Basis-family matching is the caller's responsibility.}
  \texttt{make\_auxmol} implements a priority list, not a
  family-aware selector.  A cc-pVDZ orbital basis should be paired with
  cc-pVDZ-RI explicitly.
\item \emph{The truncated metric has a floor.}  \texttt{rcond} defaults
  to $10^{-12}$ and no direction was dropped in the systems measured
  here, but the Th fallback metric is already at the $10^{9}$ condition
  number where the truncation choice starts to matter.
\end{itemize}

\section{A standalone example}

\texttt{examples/ri\_df.py} in the repository runs the H$_2$O/6-31G
CAS(2,2) case end to end: assembly against the pyscf DF solve, RI-CASCI,
the three-center and ERI derivatives against central differences, and
the JAX rule.  The core of the script is:

\begin{lstlisting}[language=Python]
import numpy as np
from pyscf import gto, scf
from diffmrc.core import ri
from diffmrc.core.casci import cs_casci

mol = gto.M(atom='O 0 0 0; H 0.7557 0 0.5880; H -0.7557 0 0.5880',
            basis='6-31g', unit='Angstrom', verbose=0)
mf = scf.RHF(mol).run(verbose=0)
auxmol, label = ri.make_auxmol(mol)
factors = ri.ri_factors(mol, auxmol)

eri_ri = ri.ri_eri_ao(mol, auxmol, factors=factors)
eri_dense = mol.intor('int2e', aosym='s1')

c = mf.mo_coeff
h1_mo, eri_mo_dense = mo_integrals(c, mf.get_hcore(), eri_dense)
b_mo, eri_mo_ri = ri.ri_eri_mo(c, factors.b, factors.jhalf)
e_ri, ci_ri = cs_casci(h1_mo[:6, :6], eri_mo_ri[:6, :6, :6, :6],
                       ncore=4, nact=2, nele_act=2, return_ci=True)

for atom in range(mol.natm):
    d3 = ri.ri_three_center_deriv(mol, auxmol, atom)
    # compare d3 against a central difference of ri_three_center
\end{lstlisting}

The measured run on the H$_2$O/6-31G system prints an assembly identity
residual of $1.1\times10^{-12}$, a dense-vs-RI RMS of
$8.3\times10^{-6}$, an RI-CASCI shift of $+1.05\times10^{-5}$~Eh, a
worst three-center derivative deviation of $3.6\times10^{-9}$, and a
worst ERI derivative deviation of $3.7\times10^{-8}$, in 6.4~s of
script time.  The companion \texttt{examples/ri\_nevpt2.py} reports the
RI-NEVPT2 correction, $+1.7\times10^{-5}$~Ha from the dense value on the
same system.
